% !TeX root = ../main.tex
\section{Introduction}
\label{sec:introduction}

Memes combine images and short text to convey meaning through cultural
references, visual symbolism, sarcasm, and other forms of multimodal
composition.  The same properties make harmful memes difficult to detect:
an image or sentence may appear benign in isolation while their combination
targets a person or group, expresses an abusive intent, or invokes a
harmful rhetorical strategy.  Effective moderation therefore requires more
than recognizing unimodal surface cues.

Prior work has addressed binary harmful or hateful meme detection with
multimodal representations, cross-modal interaction, and
contrastive learning~\cite{kiela2020hateful,kumar2022hateclipper,shah2024memeclip,mei2024rgcl}.
More recent methods have begun to predict targets or attack types and to
generate natural-language rationales~\cite{pramanick2021harmeme,sharma2022disarm,lin2024explainhm}.
These advances
move toward richer meme understanding, but fine-grained prediction,
evidence grounding, and rationale generation are still often treated as
separate or method-specific outputs.  We instead formulate harmful meme
understanding as a structured interpretation task that jointly predicts
harmfulness, target presence and granularity, primary intent, rhetorical
tactics, and multimodal relation.

A second challenge is that the meaning of a meme may depend on social,
cultural, or event-specific context that is not directly observable in its
pixels or OCR text.  Retrieval and multimodal language models can provide
such context, but acquired information is not automatically reliable
evidence.  Retrieved examples may be only superficially similar, generated
hypotheses may be unsupported, and context relevant to identifying a target
may not support the same conclusion about intent or rhetorical strategy.
This motivates our central design question: \emph{how should external
knowledge be checked against the meme before it is allowed to influence a
structured prediction?}

We propose an evidence-grounded framework that explicitly separates
knowledge acquisition from evidence verification.  Internal evidence is
extracted from global and local visual content, OCR text, and image--text
interaction, while external candidates are acquired when additional context
is useful.  Before fusion, an evidence-aware verifier evaluates each
candidate against the internal meme evidence for relevance, validity, and
task-specific support.  Only verified knowledge enters task-aware
reasoning; accepted and rejected candidates retain their provenance for
analysis.  The final model outputs the structured predictions together with
selected evidence and a concise rationale.

We evaluate cross-domain generalization by developing models exclusively on
HarMeme~\cite{pramanick2021harmeme,pramanick2021momenta} and testing them on
the held-out Facebook Hateful Memes benchmark
(FHM)~\cite{kiela2020hateful}.  No FHM example is used for model selection, prompt development, or
retrieval construction.  Beyond harmfulness, schema-validated silver
annotations enable evaluation of target, intent, rhetorical tactics, and
multimodal relation under the same domain shift.

Our contributions are threefold:
\begin{itemize}
    \item
    We formulate harmful meme understanding as structured prediction over
    harmfulness, target presence and granularity, primary intent,
    rhetorical tactics, and multimodal relation.

    \item
    We introduce an evidence-aware verification stage that evaluates
    externally acquired candidates against internal meme evidence for
    relevance, validity, and task-specific support before they can
    influence downstream multimodal reasoning.

    \item
    We conduct a cross-domain evaluation from HarMeme to FHM with
    specialized harmful-meme models and general-purpose vision--language
    baselines, and analyze structured transfer, component contributions,
    and external-knowledge filtering behavior.
\end{itemize}
